\kchapter{Services and Cluster DNS}{svc}

\begin{objectives}
\begin{itemize}
  \item the Kubernetes networking model in four rules;
  \item why Pods need a Service in front of them;
  \item ClusterIP, NodePort, LoadBalancer, ExternalName, and headless Services;
  \item \kc{port}, \kc{targetPort}, and \kc{nodePort} --- finally explained;
  \item how DNS names are built and how to debug them.
\end{itemize}
\end{objectives}

\section{The networking model}

Kubernetes networking rests on a few simple rules, implemented by your CNI plugin:

\begin{enumerate}
  \item Every Pod gets its \textbf{own IP address}.
  \item Every Pod can reach every other Pod \textbf{directly}, across nodes, without NAT.
  \item Agents on a node (like the kubelet) can reach all Pods on that node.
  \item Containers in the same Pod share the IP and talk via \kc{localhost}.
\end{enumerate}

This flat network is simple, but Pod IPs are \emph{temporary}. Every
replacement Pod gets a new IP, and a Deployment may have many Pods. Clients need
one stable address that always reaches a healthy Pod. That is a \kterm{Service}.

\section{What a Service is}

A Service gives a set of Pods (chosen by a label selector):

\begin{itemize}
  \item a stable \textbf{virtual IP} (the ClusterIP) that never changes;
  \item a stable \textbf{DNS name}, e.g.\ \kc{web.shop.svc.cluster.local};
  \item \textbf{load balancing} across all \emph{Ready} Pods behind it.
\end{itemize}

\begin{center}
\resizebox{0.95\linewidth}{!}{%
\begin{tikzpicture}[
  pod/.style={fill=k8sblue!12, draw=k8sblue, rounded corners=4pt, minimum width=2.6cm,
              minimum height=1.05cm, font=\scriptsize, align=center},
  >={Stealth[length=2.6mm]}]
  \node[fill=k8sgray!15, draw=k8sgray, rounded corners=4pt, minimum width=2.4cm,
        minimum height=1.2cm, font=\scriptsize, align=center] (c) at (-5,0)
        {\faLaptopCode\ client Pod\\\ttfamily curl http://web};
  \node[fill=k8sblue, text=white, rounded corners=6pt, minimum width=3.4cm,
        minimum height=1.5cm, align=center, font=\small\bfseries] (s) at (0,0)
        {\kimg[0.7cm]{svc}\\Service \ttfamily web\\[-1pt]\scriptsize\mdseries ClusterIP 10.96.12.7:80};
  \node[pod] (p1) at (5,1.6)  {\textbf{Pod} 10.244.1.5\\:8080 \textcolor{k8sgreen}{\faCheckCircle\ Ready}};
  \node[pod] (p2) at (5,0)    {\textbf{Pod} 10.244.2.9\\:8080 \textcolor{k8sgreen}{\faCheckCircle\ Ready}};
  \node[pod, draw=k8sred, fill=k8sred!8] (p3) at (5,-1.6) {\textbf{Pod} 10.244.2.11\\:8080 \textcolor{k8sred}{\faTimesCircle\ Not ready}};
  \draw[->, line width=1.3pt, k8sorange] (c) -- node[above, font=\scriptsize]{port 80} (s);
  \draw[->, thick, k8sblue] (s) -- node[above, sloped, font=\scriptsize]{targetPort 8080} (p1);
  \draw[->, thick, k8sblue] (s) -- (p2);
  \draw[->, thick, k8sgray, dashed] (s) -- node[below, sloped, font=\scriptsize]{excluded} (p3);
\end{tikzpicture}}
\end{center}

Behind the scenes, the EndpointSlice controller keeps a list of the IPs of all
Ready Pods matching the selector, and kube-proxy on every node programs rules so
that packets sent to the ClusterIP are forwarded to one of them.

\begin{yamlcode}{web-service.yaml}
apiVersion: v1
kind: Service
metadata:
  name: web
spec:
  type: ClusterIP             # the default
  selector:
    app: web                  # traffic goes to Ready Pods with this label
  ports:
    - name: http
      port: 80                # the port clients use on the Service
      targetPort: 8080        # the port the container listens on (or a port NAME)
      protocol: TCP
\end{yamlcode}

\begin{bashcode}{terminal}
kubectl apply -f web-service.yaml
kubectl get svc web
kubectl get endpointslices -l kubernetes.io/service-name=web
kubectl run tmp --rm -it --image=busybox:1.36 -- wget -qO- http://web
\end{bashcode}

\begin{protip}
Name your container ports (\kc{name: http}) and use \kc{targetPort: http} in the
Service. Then you can change the container port later without touching the Service.
\end{protip}

\section{Service types}

\begin{center}
\resizebox{\linewidth}{!}{%
\begin{tikzpicture}[
  lay/.style={rounded corners=6pt, minimum height=1.0cm, font=\small\bfseries,
              text=white, align=center}]
  \node[lay, fill=k8sorange, minimum width=13.6cm] at (0,2.4) {LoadBalancer\quad\mdseries\footnotesize cloud LB with public IP \faArrowRight\ NodePort \faArrowRight\ ClusterIP};
  \node[lay, fill=k8sdeep, minimum width=10.4cm] at (0,1.2) {NodePort\quad\mdseries\footnotesize every node's IP:30000-32767 \faArrowRight\ ClusterIP};
  \node[lay, fill=k8sblue, minimum width=7.2cm] at (0,0) {ClusterIP\quad\mdseries\footnotesize internal virtual IP \faArrowRight\ Pods};
  \node[font=\scriptsize\itshape, text=k8sgray] at (0,-0.9) {each type builds on the one below it};
\end{tikzpicture}}
\end{center}

\begin{center}
\ktablestart
\begin{tabularx}{\linewidth}{P{2.7cm} L L}
\kheader \kth{Type} & \kth{Reachable from} & \kth{Typical use} \\
\kc{ClusterIP} & Inside the cluster only. & Service-to-service traffic. The default and the most common. \\
\kc{NodePort} & \kc{<any-node-IP>:<nodePort>} (30000--32767). & Quick demos, bare-metal setups behind your own LB. \\
\kc{LoadBalancer} & A cloud load balancer with its own IP. & Exposing a single service publicly on a cloud. \\
\kc{ExternalName} & DNS alias (CNAME) to an outside name. & Pointing \kc{db} at \kc{mydb.rds.amazonaws.com}. \\
Headless (\kc{clusterIP: None}) & DNS returns each Pod IP. & StatefulSets, client-side load balancing. \\
\end{tabularx}
\end{center}

\begin{yamlcode}{nodeport-and-lb.yaml}
apiVersion: v1
kind: Service
metadata:
  name: web-nodeport
spec:
  type: NodePort
  selector: { app: web }
  ports:
    - port: 80
      targetPort: 8080
      nodePort: 30080         # optional; picked automatically if omitted
---
apiVersion: v1
kind: Service
metadata:
  name: web-public
spec:
  type: LoadBalancer
  externalTrafficPolicy: Local   # keep client source IP; only nodes with Pods receive traffic
  selector: { app: web }
  ports:
    - port: 80
      targetPort: 8080
---
apiVersion: v1
kind: Service
metadata:
  name: payments-db
spec:
  type: ExternalName
  externalName: payments.abc123.eu-west-1.rds.amazonaws.com
\end{yamlcode}

\begin{warningbox}
One \kc{LoadBalancer} Service = one cloud load balancer = one monthly bill. For
many HTTP services, put a single Gateway or Ingress controller behind one load
balancer and route by host/path instead (next chapter).
\end{warningbox}

\section{Cluster DNS}

CoreDNS gives every Service a name. The full form is:

\begin{center}
\colorbox{k8slight}{\large\ttfamily
\textcolor{k8sblue}{\bfseries web}.\textcolor{k8sdeep}{\bfseries shop}.\textcolor{k8snavy}{svc}.\textcolor{k8sgray}{cluster.local}}\\[4pt]
\footnotesize service \quad.\quad namespace \quad.\quad type \quad.\quad cluster domain
\end{center}

\begin{itemize}
  \item From a Pod in the \emph{same} namespace: just \kc{web}.
  \item From another namespace: \kc{web.shop} (or the full name).
  \item Headless Services in a StatefulSet also give each Pod a name:
        \kc{db-0.db.shop.svc.cluster.local}.
  \item Named ports get SRV records: \kc{\_http.\_tcp.web.shop.svc.cluster.local}.
\end{itemize}

\begin{bashcode}{terminal}
kubectl run dns --rm -it --image=busybox:1.36 -- sh
/ # nslookup web
/ # nslookup web.shop.svc.cluster.local
/ # cat /etc/resolv.conf        # search domains that make short names work
\end{bashcode}

\section{Useful Service options}

\begin{itemize}
  \item \kc{sessionAffinity: ClientIP} --- send a given client to the same Pod.
  \item \kc{internalTrafficPolicy: Local} --- only route to Pods on the caller's node.
  \item \kc{trafficDistribution: PreferClose} --- prefer topologically close endpoints (same zone).
  \item Services without a selector + manual EndpointSlices --- point a
        Service at IPs outside the cluster.
\end{itemize}

\section{Port-forward vs.\ Service}

\kc{kubectl port-forward svc/web 8080:80} tunnels from your laptop to the
cluster through the API server. It is perfect for development and debugging,
but it is \emph{not} a way to serve real users.

\begin{recap}
\begin{itemize}
  \item Pods get IPs, but IPs change; Services provide a stable IP + DNS name.
  \item Services select \emph{Ready} Pods by label and load-balance across them.
  \item \kc{port} is what clients call; \kc{targetPort} is where the container listens.
  \item ClusterIP inside; NodePort/LoadBalancer outside; headless for per-Pod DNS.
  \item DNS: \kc{service.namespace.svc.cluster.local}.
\end{itemize}
\end{recap}
